\subsection{指数映射}

命题 3. 设 \(\mathbf{G}\) 是李群。存在 \(\phi\) 的指数映射 \(\mathbf{G}\)，具有如下性质：

\begin{enumerate}
\def\labelenumi{(\roman{enumi})}
\item
  \(\phi\) 定义在加法群 L(G) 的开子群 U 上；
\item
  \(\phi (\mathbf{U})\) 是 \(\mathbf{G}\) 的开子群，且 \(\phi\) 是从解析流形 \(\mathbf{U}\) 到解析流形 \(\phi (\mathbf{U})\) 的同构；
\item
  对所有 \(\phi(nx) = \phi(x)^n\) 和所有 \(x \in \mathbf{U}\)，都有 \(n \in \mathbf{Z}\)。
\end{enumerate}

给 \(\dot{\mathrm{L}}(\mathrm{G})\) 赋予一个与其拓扑相容的范数，并使得对 \(\| [x,y]\| \leqslant \| x\| \| y\|\) 中的 \(x,y\) 有 \(\mathrm{L}(\mathrm{G})\)。令 \(\mathbf{G}_1\) 为由 \(\mathrm{L}(\mathrm{G})\) 定义的李群。令 \(\psi = \mathrm{Id}_{\mathrm{G}_1}\)，它是 \(\mathbf{G}_1\) 的一个指数映射。对所有 \(\mu >0\)，令 \(\mathrm{L}_{\mu}\) 为满足 \(x\in \mathrm{L}(\mathrm{G})\) 的 \(\| x\| < \mu\) 的集合。于是，当 \(\mu\) 充分小时，\(\mathrm{L}_{\mu}\) 是加法群 \(\mathrm{L}(\mathrm{G}),\psi (\mathrm{L}_{\mu})\)\(\mathrm{G}_1\) 的开子群，\(\psi |\mathrm{L}_{\mu}\) 是 \(\mathrm{L}_{\mu}\) 的开子群（§ 4，第 2 号，引理 3），\(\psi (\mathrm{L}_{\mu})\) 是从解析流形 \(\psi (nx) = \psi (x)^n\) 到 \(x\in \mathrm{L}_{\mu}\) 的解析流形同构，并且对所有 \(n\in \mathbf{Z}\) 和所有 \(\mathrm{L}_{\mu}\)，有 \(\mathrm{L}(\mathrm{G})\)。\(\mu\) 构成  中 0 的邻域基。由定理 1，存在 \(\mathbf{G}'\) 以及 \(\mathbf{G}\) 的开子群 \(\psi (\mathrm{L}_{\mu})\)，使得  与 \(\mathbf{G}'\) 同构，从而得到命题。

命题 4. 设 G 是李群，\(\phi\) 是 G 的单射指数映射。假设 \(p > 0\)。对 L(G) 中所有 \(x, y\)，

\begin{enumerate}
\def\labelenumi{(\arabic{enumi})}
\tightlist
\item
\end{enumerate}

\[
x + y = \lim _ {n \rightarrow + \infty} p ^ {- n} \phi^ {- 1} (\phi (p ^ {n} x) \phi (p ^ {n} y))\tag{1}
\]

\[
[ x, y ] = \lim _ {n \rightarrow + \infty} p ^ {- 2 n} \phi^ {- 1} (\phi (p ^ {n} x) \phi (p ^ {n} y) \phi (- p ^ {n} x) \phi (- p ^ {n} y)).\tag{2}
\]

这些是 § 4，第 3 号，命题 4 的特例。

